% Ajout de réactifs à un système à l'équilibre : ions fer(III) et thiocyanate (exercice 4, chapitre 7)
\begin{tikzpicture}[x=1cm,y=1cm]
\colorlet{rougeA}{orange!55!red!45!white}
\colorlet{rougeB}{red!60!black!75}
\foreach \x/\c/\t in {0/rougeA/{tube 1 :\\témoin},2.6/rougeB/{tube 2 : ajout\\d'ions \ce{Fe^3+}},5.2/rougeB/{tube 3 : ajout\\d'ions \ce{SCN-}}}{%
  \pic at (\x,0) {tube={hauteur=2.4,niveau=1.5,liquide=\c}};
  \node[pcLegende,align=center,anchor=north] at (\x,-0.15) {\t};}
\end{tikzpicture}
\par\smallskip\legende{Les trois tubes contiennent initialement le même mélange à l'équilibre, de couleur rouge orangé.}
